%% chapters/06-kesimpulan.tex
%% Bab VI — Kesimpulan dan Saran
%% Aturan: lihat .cursor/rules/11-bab6-kesimpulan.mdc

\chapter{Kesimpulan dan Saran}
\label{bab:kesimpulan}

\section{Kesimpulan}
\label{sec:bab6_kesimpulan}

Disertasi ini telah meneliti prediksi sisa umur pakai (RUL) bantalan
gelinding menggunakan arsitektur hibrida \emph{Mamba-xLSTM-Net} dengan
modul interpretabilitas berbasis \emph{Top-k Sparse Autoencoder}. Tiga
pertanyaan penelitian yang dirumuskan pada
Subbab~\ref{subsec:bab1_masalah} dijawab sebagai berikut.

\subsection{Jawaban atas Pertanyaan Penelitian}
\label{subsec:bab6_jawaban_rq}

\paragraph{Pertanyaan penelitian pertama (RQ1)}
mempertanyakan apakah kombinasi \emph{state-space model} selektif
(Mamba) dengan \emph{matrix-memory recurrent} (mLSTM) memberikan
kinerja prediksi RUL yang lebih baik daripada masing-masing komponen
secara terpisah maupun arsitektur Transformer. Hasil eksperimen
(Subbab~\ref{subsec:bab5_main_xjtusy} dan
Subbab~\ref{subsec:bab5_main_phm2012}) menunjukkan penurunan RMSE
sebesar XX~\% pada XJTU-SY dan XX~\% pada PHM2012 dibanding baseline
terkuat. RQ1 dijawab \emph{ya}, dengan margin signifikan secara
statistik. Hasil ini mendukung Hipotesis~H1.

\paragraph{Pertanyaan penelitian kedua (RQ2)}
% Jawaban RQ2: skema interpretabilitas SAE + pemetaan ke BPFO/BPFI/BSF.
\dots

\paragraph{Pertanyaan penelitian ketiga (RQ3)}
% Jawaban RQ3: generalisasi lintas dataset.
\dots

\subsection{Kontribusi Disertasi}
\label{subsec:bab6_kontribusi}

Disertasi ini menyumbangkan empat kontribusi utama bagi bidang
prognostik bantalan gelinding:

\begin{enumerate}
  \item Mengusulkan arsitektur hibrida \emph{Mamba-xLSTM-Net} yang
        mengintegrasikan \emph{state-space model} selektif dan
        \emph{matrix-memory recurrent} dengan mekanisme fusi yang
        diturunkan dari sifat non-stasioner sinyal getaran.
  \item Memformulasikan modul interpretabilitas \emph{post-hoc}
        berbasis \emph{Top-k Sparse Autoencoder} yang bekerja
        terhadap representasi laten model RUL tanpa mendegradasi
        akurasi prediksi.
  \item Membangun korespondensi empiris antara fitur laten SAE dan
        frekuensi karakteristik bantalan (BPFO, BPFI, BSF, FTF) yang
        memberikan bukti bahwa model belajar representasi konsisten
        dengan teori klasik diagnosis getaran.
  \item Menyediakan \emph{pipeline reproducible} berikut studi ablasi
        komprehensif pada dua dataset publik (PHM2012 dan XJTU-SY)
        dengan baseline yang luas.
\end{enumerate}

\subsection{Implikasi}
\label{subsec:bab6_implikasi}

% Implikasi keilmuan: kontribusi terhadap literatur RUL dan deep learning.
% Implikasi praktis: penerapan untuk industri pemeliharaan.
% Implikasi metodologis: pelajaran tentang menggabungkan SSM dengan recurrent.
\dots

\section{Saran Pekerjaan Lanjut}
\label{sec:bab6_saran}

\subsection{Pengembangan Metodologis}
\label{subsec:bab6_saran_metode}

\paragraph{Validasi kausal terhadap konsep SAE.}
Klaim interpretabilitas dalam disertasi ini didasarkan pada korelasi
statistik antara aktivasi SAE dan komponen spektral. Pekerjaan lanjut
yang menjanjikan adalah validasi \emph{causal} melalui intervensi:
secara sengaja menambahkan komponen frekuensi karakteristik buatan ke
sinyal, lalu memeriksa apakah fitur SAE yang relevan menyala.

\paragraph{Kuantifikasi ketidakpastian prediksi.}
% Bayesian deep learning, MC Dropout, deep ensembles untuk RUL.
\dots

\subsection{Perluasan Domain Aplikasi}
\label{subsec:bab6_saran_domain}

\paragraph{Pengujian pada bantalan jenis lain.}
% Spherical roller, tapered roller, magnetic bearing.
\dots

\paragraph{Pengujian pada kondisi industri sebenarnya.}
% Data noise tinggi, beban variabel, kondisi tidak terkontrol.
\dots

\subsection{Aspek Implementasi}
\label{subsec:bab6_saran_impl}

\paragraph{Kompresi model untuk \emph{edge deployment}.}
% Pruning, quantization untuk on-device inference real-time.
\dots

\paragraph{Integrasi dengan sistem PdM industri.}
% PLC, SCADA, MQTT pipelines.
\dots
